\documentclass[12pt]{article}
\setlength{\oddsidemargin}{0in}
\setlength{\evensidemargin}{0in}
\setlength{\textwidth}{6.5in}
\setlength{\parindent}{0in}
\setlength{\parskip}{\baselineskip}

\usepackage{amsmath,amsfonts,amssymb}

\title{Math 291 Challenge Set}

\begin{document}

Math 291 Fall 2025\hfill{} Problem Set 1 \\
Jeffrey Lu, Robert, Matthew, Luke

\hrulefill{}

DATA: Consider \(\mathbf{v}_1 = (1, -1, 0, 1, 1), \mathbf{v}_2 = (0, 1, 1, 0, 1), \mathbf{v}_3 = (-1, 3, 2, -1, 1), \mathbf{v}_4 = (1, 0, 1, 1, 2), \mathbf{v}_5 = (0, 0, 0, 1, 0), \mathbf{v}_6 = (0, 0, 1, 0, 1), \mathbf{v}_7 = (0, 0, 0, 0, 1)\). For \(1 \leq i \leq j \leq 7\), let \(A^{(i,j)}\) denote the span of \(\mathbf{v}_i, \ldots, \mathbf{v}_j\).
\begin{enumerate}
    \item{} \textit{(a) Perform Gram-Schmidt on the input sequence of seven vectors \(\mathbf{v_1} , \ldots, \mathbf{v_7}\) to obtain the output sequence \(\mathbf{u_1}, \ldots, \mathbf{u_r}\). Do by hand first, and check your answer in other ways.
              (b) Write each input vector as a linear combination of output vectors.
              (c) Do you expect the same output sequence if the input sequence is reversed? What aspect of the output sequence will not change?}

          \begin{enumerate}
              \item{}

                    Choose the first vector to be pivotal and normalize:
                    \[
                        \mathbf{u}_1=\frac{\mathbf{v}_1}{\|\mathbf{v}_1\|}=\tfrac{1}{2}(1,-1,0,1,1)
                    \]

                    Next, apply Gram–Schmidt to \(\mathbf{v}_2\):
                    \begin{align*}
                        \mathbf{w}_2 & = \mathbf{v}_2- (\mathbf{u}_1\cdot\mathbf{v}_2)\mathbf{u}_1 \\
                        \mathbf{w}_2 & = \mathbf{v}_2-\mathbf{0} = (0,1,1,0,1)
                    \end{align*}
                    Normalizing \(\mathbf{w}_2\):
                    \[
                        \mathbf{u}_2 = \frac{\mathbf{w}_2}{\|\mathbf{w}_2\|}=\tfrac{1}{\sqrt{3}}(0,1,1,0,1)
                    \]
                    For \(\mathbf{v}_3\):
                    \begin{align*}
                        \mathbf{w}_3 & = \mathbf{v}_3-(\mathbf{u}_1\cdot\mathbf{v}_3)\mathbf{u}_1- (\mathbf{u}_2\cdot\mathbf{v}_3)\mathbf{u}_2 \\
                                     & = \mathbf{v}_3+\mathbf{v}_1-2\mathbf{v}_2= (0,0,0,0,0)                                                  \\
                                     & = \mathbf{0}
                    \end{align*}
                    Thus \(\mathbf{v}_3\) is not a pivotal vector. \\

                    For \(\mathbf{v}_4\):
                    \begin{align*}
                        \mathbf{w}_4 & = \mathbf{v}_4- (\mathbf{u}_1\cdot\mathbf{v}_4)\mathbf{u}_1-(\mathbf{u}_2\cdot\mathbf{v}_4)\mathbf{u}_2 \\
                                     & = \mathbf{v}_4-\mathbf{v}_1-\mathbf{v}_2= (0,0,0,0,0)                                                   \\
                                     & = \mathbf{0}
                    \end{align*}
                    Thus \(\mathbf{v}_4\) is also not a pivotal vector \\

                    For \(\mathbf{v}_5\):
                    \begin{align*}
                        \mathbf{w}_5 & =\mathbf{v}_5-(\mathbf{u}_1\cdot\mathbf{v}_5)\mathbf{u}_1-(\mathbf{u}_2\cdot\mathbf{v}_5)\mathbf{u}_2 \\
                                     & =\mathbf{v}_5-\tfrac{1}{4}\mathbf{v_1}-\mathbf{0}                                                     \\
                                     & =\tfrac{1}{4} (-1,1,0,3,-1)
                    \end{align*}
                    Normalizing, we get our third orthonormal vector
                    \[
                        \mathbf{u}_3=\frac{\mathbf{w}_5}{\|\mathbf{w}_5\|}
                        =\tfrac{1}{\sqrt{12}}(-1,1,0,3,-1)
                    \]

                    For \(\mathbf{v}_6\):
                    \begin{align*}
                        \mathbf{w}_6 & =\mathbf{v}_6-(\mathbf{u}_1\cdot\mathbf{v}_6)\mathbf{u}_1-(\mathbf{u}_2\cdot\mathbf{v}_6)\mathbf{u}_2-(\mathbf{u}_3\cdot\mathbf{v}_6)\mathbf{u}_3 \\
                                     & =\mathbf{v}_6-\tfrac{1}{4}\mathbf{v_1}-\tfrac{2}{3}\mathbf{v_2}+\tfrac{1}{12} (-1,1,0,3,-1)                                                       \\
                                     & =\tfrac{1}{3} (-1,-1,1,0,0)
                    \end{align*}
                    Normalizing, we get our fourth orthonormal vector
                    \[
                        \mathbf{u}_4=\frac{\mathbf{w}_6}{\|\mathbf{w}_6\|} =\tfrac{1}{\sqrt{3}}(-1,-1,1,0,0)
                    \]

                    Finally, for \(\mathbf{v}_7\):
                    \begin{align*}
                        \mathbf{w}_7 & =\mathbf{v}_7-(\mathbf{u}_1\cdot\mathbf{v}_7)\mathbf{u}_1-(\mathbf{u}_2\cdot\mathbf{v}_7)\mathbf{u}_2-(\mathbf{u}_3\cdot\mathbf{v}_7)\mathbf{u}_3-(\mathbf{u}_4\cdot\mathbf{v}_7)\mathbf{u}_4 \\
                                     & =\mathbf{v}_7-\tfrac{1}{4}\mathbf{v}_1-\tfrac{1}{3}\mathbf{v}_2+\tfrac{1}{12} (-1,1,0,3,-1)-\mathbf{0}                                                                                        \\
                                     & =\tfrac{1}{3} (-1,0,-1,0,1)
                    \end{align*}
                    Normalizing, we get our fifth and final orthonormal vector, completing the basis for \(\mathbb{R}^5\):
                    \[
                        \mathbf{u}_5=\frac{                                                                                                                                                                                                                                                                                                                                                                                                                                                                                                                                                                                   \mathbf{w}_7}{\|\mathbf{w}_7\|} =\tfrac{1}{\sqrt{3}}(-1,0,-1,0,1)
                    \]

              \item{}
                    Using the fundamental theorem of orthogonal vectors and the fact that our output vectors are orthonormal,
                    \[
                        \mathbf{v}_j = \sum_{i=1}^{r} (\mathbf{u}_i \cdot \mathbf{v}_j) \mathbf{u}_i \quad \text{for } j=1,\ldots,7
                    \]
                    We can rearrange the equations used in part (a) to get input vectors in terms of the output vectors.
                    \begin{align*}
                        \mathbf{v_1} = \|\mathbf{v_1}\|\mathbf{u}_1 = 2\mathbf{u}_1 \\
                        \mathbf{v_2}  = \|\mathbf{w_2}\|\mathbf{u}_2 = \sqrt{3}\mathbf{u}_2
                    \end{align*}
                    We already found that \(\mathbf{v}_3\) is a linear combination of \(\mathbf{v}_1\) and \(\mathbf{v}_2\):
                    \begin{align*}
                        \mathbf{v_3} & = -\mathbf{v_1} + 2\mathbf{v_2}          \\
                                     & = -2\mathbf{u}_1 + 2\sqrt{3}\mathbf{u}_2
                    \end{align*}
                    Again, \(\mathbf{v}_4\) is also a linear combination of \(\mathbf{v}_1\) and \(\mathbf{v}_2\):
                    \begin{align*}
                        \mathbf{v_4} & = \mathbf{v_1} + \mathbf{v_2}          \\
                                     & = 2\mathbf{u}_1 + \sqrt{3}\mathbf{u}_2
                    \end{align*}
                    The remaining vectors are the same process:
                    \begin{align*}
                        \mathbf{v_5} & = \mathbf{w_5} + \tfrac{1}{4}\mathbf{v_1} + 0\mathbf{u}_2                                                                              \\
                                     & = \|\mathbf{w_5}\|\mathbf{u}_3 + \tfrac{1}{4}\mathbf{v_1}                                                                              \\
                                     & = \tfrac{\sqrt{12}}{4}\mathbf{u}_3 + \tfrac{1}{4}(2\mathbf{u}_1)                                                                       \\
                                     & = \tfrac{1}{2}\mathbf{u}_1 + \tfrac{\sqrt{3}}{2}\mathbf{u}_3                                                                           \\
                        \\
                        \mathbf{v_6} & = \mathbf{w_6} + \tfrac{1}{4}\mathbf{v_1} + \tfrac{2}{3}\mathbf{v_2} - \tfrac{1}{\sqrt{12}}\mathbf{u}_3                                \\
                                     & = \|\mathbf{w_6}\|\mathbf{u}_4 + \tfrac{1}{4}\mathbf{v_1} + \tfrac{2}{3}\mathbf{v_2} - \tfrac{\sqrt{3}}{6}\mathbf{u}_3                 \\
                                     & = \tfrac{\sqrt{3}}{3}\mathbf{u}_4 + \tfrac{1}{4}(2\mathbf{u}_1) + \tfrac{2}{3}(\sqrt{3}\mathbf{u}_2) - \tfrac{\sqrt{3}}{6}\mathbf{u}_3 \\
                                     & = \tfrac{1}{2}\mathbf{u}_1 + \tfrac{2\sqrt{3}}{3}\mathbf{u}_2 - \tfrac{\sqrt{3}}{6}\mathbf{u}_3 + \tfrac{\sqrt{3}}{3}\mathbf{u}_4      \\
                        \\
                        \mathbf{v_7} & = \mathbf{w_7} + \tfrac{1}{4}\mathbf{v_1} + \tfrac{1}{3}\mathbf{v_2} - \tfrac{1}{\sqrt{12}}\mathbf{u}_3 + 0\mathbf{u}_4                \\
                                     & = \|\mathbf{w_7}\|\mathbf{u}_5 + \tfrac{1}{4}\mathbf{v_1} + \tfrac{1}{3}\mathbf{v_2} - \tfrac{\sqrt{3}}{6}\mathbf{u}_3                 \\
                                     & = \tfrac{\sqrt{3}}{3}\mathbf{u}_5 + \tfrac{1}{4}(2\mathbf{u}_1) + \tfrac{1}{3}(\sqrt{3}\mathbf{u}_2) - \tfrac{\sqrt{3}}{6}\mathbf{u}_3 \\
                                     & = \tfrac{1}{2}\mathbf{u}_1 + \tfrac{\sqrt{3}}{3}\mathbf{u}_2 - \tfrac{\sqrt{3}}{6}\mathbf{u}_3 + \tfrac{\sqrt{3}}{3}\mathbf{u}_5
                    \end{align*}
                    Collecting all the equations, we have:
                    \[
                        \begin{array}{rcl}
                            \mathbf{v}_1 & = & 2\mathbf{u}_1 + 0\mathbf{u}_2 + 0\mathbf{u}_3 + 0\mathbf{u}_4 + 0\mathbf{u}_5,                                                                   \\[6pt]
                            \mathbf{v}_2 & = & 0\mathbf{u}_1 + \sqrt{3}\mathbf{u}_2 + 0\mathbf{u}_3 + 0\mathbf{u}_4 + 0\mathbf{u}_5,                                                            \\[6pt]
                            \mathbf{v}_3 & = & -2\mathbf{u}_1 + 2\sqrt{3}\mathbf{u}_2 + 0\mathbf{u}_3 + 0\mathbf{u}_4 + 0\mathbf{u}_5,                                                          \\[6pt]
                            \mathbf{v}_4 & = & 2\mathbf{u}_1 + \sqrt{3}\mathbf{u}_2 + 0\mathbf{u}_3 + 0\mathbf{u}_4 + 0\mathbf{u}_5,                                                            \\[6pt]
                            \mathbf{v}_5 & = & \tfrac{1}{2}\mathbf{u}_1 + 0\mathbf{u}_2 + \tfrac{\sqrt{3}}{2}\mathbf{u}_3  + 0\mathbf{u}_4 + 0\mathbf{u}_5,                                     \\[6pt]
                            \mathbf{v}_6 & = & \tfrac{1}{2}\mathbf{u}_1 + \tfrac{2\sqrt{3}}{3}\mathbf{u}_2 - \tfrac{\sqrt{3}}{6}\mathbf{u}_3 + \tfrac{\sqrt{3}}{3}\mathbf{u}_4 + 0\mathbf{u}_5, \\[6pt]
                            \mathbf{v}_7 & = & \tfrac{1}{2}\mathbf{u}_1 + \tfrac{\sqrt{3}}{3}\mathbf{u}_2 - \tfrac{\sqrt{3}}{6}\mathbf{u}_3 + 0\mathbf{u}_4 + \tfrac{\sqrt{3}}{3}\mathbf{u}_5.
                        \end{array}
                    \]


              \item{}
                    If the input sequence is reversed, we do not expect the same output sequence.
                    This is because if Gran-Schmidt is applied to a set a vectors in a different order, you may get a different set of pivotal vectors, and therefore the orthonormal basis that you end up with will also be different.
                    However, the span of the orthonormal basis will always be the same, in this case spanning all of $\mathbb{R}^5$, no matter the order you applied Gram–Schmidt.

          \end{enumerate}

          \vspace{1em}

    \item{}  Recall that a subspace S of $\mathbb{R}^n$ is a subset of $\mathbb{R}^n$ satisfying two properties. What are they? Specify in detail - so, not just the
          name of the properties. (Notice that the span of any set is always a subspace, as it satisfies both these properties.)
          \begin{enumerate}
              \item{} Closed under vector addition: For any two vectors \(\mathbf{v}_1, \mathbf{v}_2 \in \text{set } S\), the sum \(\mathbf{v}_1 + \mathbf{v}_2\) is also in set \(S\).
              \item{} Closed under scalar multiplication: For any vector \(\mathbf{v} \in \text{set } S\), when multiplied by any scalar \(c \in \mathbb{R}\), the resulting vector \(c\mathbf{v}\) is also in set \(S\).
          \end{enumerate}

          \vspace{1em}

    \item{} Every non-zero subspace of R3 is the span of an orthonormal set of , , or vectors, which we respectively recognize
          geometrically as , , or through the origin in R3.
          \begin{enumerate}
              \item{}1,2,3

              \item{} line, plane, or entire space of $\mathbb{R}^3$
          \end{enumerate}

          \vspace{1em}

    \item{} For each subspace A(1,j) j = 1, \dots, 7 obtain
          \begin{enumerate}
              \item{}
                    We already found the orthonormal basis by perfoming Gram–Schmidt on \(\mathbf{v_1} , \ldots, \mathbf{v_7}\).
                    Decreasing j will take vectors away, shrinking the orthonormal basis, by only using  orthonormal vectors found before and from the jth vector during Gram–Schmidt.
                    \begin{align*}
                        j & =1 \rightarrow{} \{\mathbf{u}_1\}                                                           \\
                        j & =2 \rightarrow{} \{\mathbf{u}_1\, \mathbf{u}_2 \}                                           \\
                        j & =3 \rightarrow{} \{\mathbf{u}_1\, \mathbf{u}_2 \}                                           \\
                        j & =4 \rightarrow{} \{\mathbf{u}_1\, \mathbf{u}_2 \}                                           \\
                        j & =5 \rightarrow{} \{\mathbf{u}_1\, \mathbf{u}_2, \mathbf{u}_3 \}                             \\
                        j & =6 \rightarrow{} \{\mathbf{u}_1\, \mathbf{u}_2, \mathbf{u}_3, \mathbf{u}_4 \}               \\
                        j & =7 \rightarrow{} \{\mathbf{u}_1\, \mathbf{u}_2, \mathbf{u}_3, \mathbf{u}_4, \mathbf{u}_5 \}
                    \end{align*}

              \item{}
                    The dimension of the subspace is equal to the number of orthonormal vectors in the ONB.
                    \begin{align*}
                        j & =1 \rightarrow{} dim = 1 \\
                        j & =2 \rightarrow{} dim = 2 \\
                        j & =3 \rightarrow{} dim = 2 \\
                        j & =4 \rightarrow{} dim = 2 \\
                        j & =5 \rightarrow{} dim = 3 \\
                        j & =6 \rightarrow{} dim = 4 \\
                        j & =7 \rightarrow{} dim = 5
                    \end{align*}
              \item{}
                    Any vector in the subspace can be parameterized into a linear combination of a orthonormal basis of the subspace.
                    \begin{align*}
                        P (t_1, \dots, t_k) = \sum_{j=1}^{k} t_j\mathbf{u}_j
                    \end{align*}
                    This function preserves dot product as demonstrated below:
                    \begin{align*}
                        P (\mathbf{t}) \cdot{} P (\mathbf{s}) & = \sum_{j=1}^{k} t_j\mathbf{u}_j \cdot{} \sum_{i=1}^{k} s_i\mathbf{u}_i    \\
                                                              & = \sum_{i=1}^{k}\sum_{j=1}^{k} t_js_i (\mathbf{u}_j  \cdot{} \mathbf{u}_i) \\
                                                              & = \sum_{i=1}^{k} t_js_i                                                    \\
                                                              & = \mathbf{t} \cdot{} \mathbf{s}
                    \end{align*}
              \item{}
                    Since the vectors are in $\mathbb{R}^5$, the orthogonal complement of the subspace is just \(A^{(j+1,7)}\), expect for the j=7 where it's just the zero subspace.
                    \begin{align*}
                        j & =1 \rightarrow{} \{\mathbf{u}_2, \mathbf{u}_3, \mathbf{u}_4, \mathbf{u}_5\} \\
                        j & =2 \rightarrow{} \{\mathbf{u}_3, \mathbf{u}_4, \mathbf{u}_5\}               \\
                        j & =3 \rightarrow{} \{\mathbf{u}_3, \mathbf{u}_4, \mathbf{u}_5\}               \\
                        j & =4 \rightarrow{} \{\mathbf{u}_3, \mathbf{u}_4, \mathbf{u}_5\}               \\
                        j & =5 \rightarrow{} \{\mathbf{u}_4, \mathbf{u}_5\}                             \\
                        j & =6 \rightarrow{} \{\mathbf{u}_5 \}                                          \\
                        j & =7 \rightarrow{} \{0\}
                    \end{align*}

          \end{enumerate}
          \vspace{0.5em}

    \item{}
          In the next few parts, we’ll develop some machinery. You might also try the last problem now to get a feeling for the concrete details.

          The set obtained by “translating” a subspace $V$ by a point $\mathbf{p}$, so
          \[
              F = p + V := \{p + v : v \in V\},
          \]
          is called an \textbf{affine subset} of $\mathbb{R}^n$. (The $p$ is replaceable by any $p + v$ for $v \in V$ without changing the set $F$). The dimension of $V$ is called the dimension of $F$.

          \begin{enumerate}
              \item[(a)] Show that a line with parametric form $\mathbf{x}_0 + t\mathbf{a}$ in $\mathbb{R}^3$ is an affine subset of $\mathbb{R}^3$. Identify $\mathbf{p}, V$. So, is the dimension of the line?

              \item[(b)] Consider $\mathbf{n} \neq \mathbf{0}, \mathbf{x}_0 \in \mathbb{R}^n$. Let
                    \[
                        H = \{\mathbf{x} \in \mathbb{R}^n : (\mathbf{x} - \mathbf{x}_0) \cdot \mathbf{n} = 0\}.
                    \]
                    Show this is an affine subset of dimension $n-1$. Such subsets are called \textbf{hyperplanes}.
          \end{enumerate}


          \textbf{SOLUTION}:

          \begin{enumerate}
              \item{}
                    The parameterization of a line is in the same form of an affine subset in $\mathbb{R}^n$.
                    The base point $\mathbf{x}_0$ is the same as the point $\mathbf{p}$. While $V = t\mathbf{a} = span(\mathbf{a})$.
                    $\mathbf{a}$ represents the subspace V because any two vectors lying on the line added together will still lie on the line and if $\mathbf{a}$ is scaled by any real number it also still lies on the line.
                    dim(V) = dim(F) = dim($\mathbf{x}_0+t\mathbf{a}$) = 1
              \item{} Expanding H you get the hyperplane in the form $a\mathbf{x}+b\mathbf{y}+\dots+c\mathbf{z} + d$, if d=0 then the plane goes through the origin and is not a affine subset, otherwise it is.


          \end{enumerate}
    \item{}
          If $\mathbf{u}_1, \ldots, \mathbf{u}_k$ is an orthonormal (generating) spanning set (so, orthonormal basis) of a subspace V of $\mathbb{R}^n$.
          Let us denote by $Proj_V (\mathbf{x})$ or $\mathbf{x}_V $ the vector  $\sum_{i=1}^{k} (\mathbf{x} \cdot \mathbf{u}_i)\mathbf{u}_i$.
          Prove:

          \begin{enumerate}
              \item{} For \( \mathbf{x}, \mathbf{y} \in \mathbb{R}^n, (\mathbf{x} + \mathbf{y})_{V} = \mathbf{x}_{V} + \mathbf{y}_{V} \): \\

                    This can be easily proved by just distributing the dot products, vector multiplication, and seperating the summations like so:
                    \begin{align*}
                        (\mathbf{x}+\mathbf{y})_V & = \sum_{j=1}^k ((\mathbf{x}+\mathbf{y})\cdot{\mathbf{u}_j})\mathbf{u}_j                                             \\
                                                  & = \sum_{j=1}^k (\mathbf{x}\cdot{\mathbf{u}_j}+\mathbf{y}\cdot{\mathbf{u}_j})\mathbf{u}_j                            \\
                                                  & = \sum_{j=1}^k (\mathbf{x}\cdot{\mathbf{u}_j})\mathbf{u}_j+ (\mathbf{y}\cdot{\mathbf{u}_j})\mathbf{u}_j             \\
                                                  & = \sum_{j=1}^k (\mathbf{x}\cdot{\mathbf{u}_j})\mathbf{u}_j+\sum_{i=1}^k (\mathbf{y}\cdot{\mathbf{u}_i})\mathbf{u}_i \\
                                                  & = \mathbf{x}_V+\mathbf{y}_V
                    \end{align*}
              \item{} For \(\mathbf{x},\mathbf{y}\in{\mathbb{R}^n}, (\alpha\mathbf{x})_V = \alpha(\mathbf{x}_V) \): \\

                    Since \(\alpha \in \mathbb{R}\), it can be factored out like so:
                    \begin{align*}
                        (\alpha\mathbf{x})_V & = \sum_{i=1}^{k} ((\alpha\mathbf{x}) \cdot{} \mathbf{u}_i)\mathbf{u}_i \\
                                             & = \sum_{i=1}^{k} (\alpha(\mathbf{x} \cdot{} \mathbf{u}_i))\mathbf{u}_i \\
                                             & = \alpha\sum_{i=1}^{k} (\mathbf{x} \cdot{} \mathbf{u}_i)\mathbf{u}_i   \\
                                             & = \alpha(\mathbf{x}_V)                                                 \\
                    \end{align*}
              \item{} If \(\mathbf{x}\in{V}\text{, then }\mathbf{x}_V = \mathbf{x}\): \\

                    If $\mathbf{x}$ is an element of subspace V, then it can also be written as a linear combination of any orthonormal basis in V
                    \begin{align*}
                        \mathbf{x}_V & = \sum_{i=1}^{k} (\mathbf{x} \cdot{} \mathbf{u}_i)\mathbf{u}_i \\
                                     & = \mathbf{x}
                    \end{align*}

          \end{enumerate}
    \item{} \emph{As we saw in class, $\mathbf{x} = \mathbf{x}_V + (\mathbf{x} - \mathbf{x}_V)$ is the orthogonal decomposition of any vector in $\mathbb{R}^n$ (uniquely) as a sum of one in $V$ and one in $V^\perp$. Use the argument in question 5 of homework 3 to show that $\mathbf{p} + (\mathbf{x} - \mathbf{p})_V$ is the unique point in the affine subset $F = \mathbf{p} + V$ that is closest to $\mathbf{x}$ -- where the minimum of $\|\mathbf{x} - (\mathbf{p} + \mathbf{v})\|$ (or its square) among all $\mathbf{v} \in V$ is attained.} \\

          \textbf{SOLUTION:}

          To find the minimum of closest distance between the affine subset F and vector $\mathbf{x}$ take the orthogonal decomposition of $\|\mathbf{x} - (\mathbf{p} + \mathbf{v})\|$ with regard to subspace V like so:
          \begin{align*}
              \|\mathbf{x} - (\mathbf{p} + \mathbf{v})\| & = \|(\mathbf{x} - (\mathbf{p} + \mathbf{v}))_V\| + \|\mathbf{x} - (\mathbf{p} + \mathbf{v})-(\mathbf{x} - (\mathbf{p} + \mathbf{v}))_V\| \\
                                                         & =\|\mathbf{x}_V - \mathbf{p}_V - \mathbf{v}\| + \|\mathbf{x} - \mathbf{p} - \mathbf{v}-\mathbf{x}_V + \mathbf{p}_V + \mathbf{v}\|        \\
                                                         & =\|\mathbf{x}_V - \mathbf{p}_V - \mathbf{v}\| + \|\mathbf{x} - \mathbf{p} -\mathbf{x}_V + \mathbf{p}_V \|
          \end{align*}
          The vector $\mathbf{v}$ is already in subspace V, so it stays the same as proven in 6(c).
          Additonally, the orthogonal component is now independent of vector $\mathbf{v}$, therefore it can be ignored.

          Next, we minimize the component projected onto subspace V by choosing a specific $\mathbf{v}$, to find this we set the expression to 0 and solve for $\mathbf{v}$:
          \begin{align*}
              \mathbf{x}_V - \mathbf{p}_V - \mathbf{v} & = 0          \\
              (\mathbf{x} - \mathbf{p})_V              & = \mathbf{v}
          \end{align*}
          Now, find the unique point in the affine subset F by plugging in the minimizied $\mathbf{v}$ in to \(\mathbf{f}=\mathbf{p}+\mathbf{v} | (\mathbf{f} \in F)\), $\mathbf{f}$ is the unique point minimizing distance from F to $\mathbf{x}$
          \begin{align*}
              \mathbf{f} & =\mathbf{p}+\mathbf{v}                  \\
                         & =\mathbf{p}+(\mathbf{x} - \mathbf{p})_V
          \end{align*}
          This matches the unique point stated earlier.

    \item{} If $V, W$ are subspaces of same $\mathbb{R}^n$, their difference defined by $V - W = \{\mathbf{v} - \mathbf{w} : \mathbf{v} \in V, \mathbf{w} \in W\}$ is also a subspace. (Note: because $W$ is a subspace, the minus can be absorbed into the general $\mathbf{w}$ and so this subset is same as $V + W$.) Therefore, if $\mathbf{p} + V$ and $\mathbf{q} + W$ are affine subsets, so is their difference $\mathbf{p} - \mathbf{q} + V - W$.
          \begin{enumerate}
              \item{} Show that the distance between affine subsets $\mathbf{p} + V$ and $\mathbf{q} + W$ is $\|\mathbf{p} - \mathbf{q} - (\mathbf{p} - \mathbf{q})_{(V - W)}\|$. \\

                    \textbf{SOLUTION:}
                    We learned from above that the distance between two affine subsets is equivalent to the distance between a point created from the difference between the two base points and a subspace created by the combination of the two subspaces of the affine subsets V and W.

                    Therefore, we can take the orthogonal decomposition of this new expression to uncover the distance between any two affine subsets like so:
                    \begin{align*}
                        \| \mathbf{p}-\mathbf{q} + V - W \|\  & = \| (\mathbf{p}-\mathbf{q})_{V-W} + (V - W)\| +                                                            \\
                                                              & \quad{} \|(\mathbf{p}-\mathbf{q}) + (V - W) -((\mathbf{p}-\mathbf{q})+(V-W))_{V-W} \|                       \\
                                                              & = \| (\mathbf{p}-\mathbf{q})_{V-W} + (V - W)\| + \|(\mathbf{p}-\mathbf{q})-(\mathbf{p}-\mathbf{q})_{V-W} \|
                    \end{align*}
                    We can minimize the projection onto (V-W) by setting (V-W) to a unique point because \((\mathbf{p}-\mathbf{q})_(V-W)\) lies in (V-W), so there exists a vector within (V-W) that is equal to it and therefore set the entire projected component to zero.
                    Therefore the distance between two affine subsets is equal to the norm of only the component orthogonal to the subspace (V-W):
                    \[
                        \| \mathbf{p}-\mathbf{q} + V - W \| = \|(\mathbf{p}-\mathbf{q})-(\mathbf{p}-\mathbf{q})_{V-W} \| \quad \square
                    \]

              \item{} If $\mathbf{u}_1, \ldots, \mathbf{u}_l$ is an ONB of $V$ and $\mathbf{w}_1, \ldots, \mathbf{w}_m$ an ONB of $W$, would their ‘union’, $\mathbf{u}_1, \ldots, \mathbf{u}_l, \mathbf{w}_1, \ldots, \mathbf{w}_m$ span the set $V - W$? Do you expect it to be an orthonormal set? If no, propose a strategy to get an orthonormal basis for $V - W$. \\

                    \textbf{SOLUTION:}

                    Yes, the union of the two subspaces would span the set $V-W$.
                    This is because if you define \(\mathbf{x} \in{} (V-W)\), it is also in (V+W) because V+W = V-W. We also define arbritary vectors:
                    \begin{align*}
                        \mathbf{v} & = \sum_{i=1}^l(\mathbf{v}\cdot{} \mathbf{u_i})\mathbf{u_i} \in{V} \\
                        \mathbf{w} & = \sum_{j=1}^m(\mathbf{w}\cdot{} \mathbf{u_j})\mathbf{u_j} \in{W}
                    \end{align*}
                    Because \(\mathbf{x} \in{(V+W)}\), this means it can also be written as:
                    \begin{align*}
                        \mathbf{x} & = \mathbf{v} + \mathbf{w}                                                                                             \\
                                   & = \sum_{i=1}^l(\mathbf{v}\cdot{} \mathbf{u_i})\mathbf{u_i} + \sum_{j=1}^m(\mathbf{w}\cdot{} \mathbf{u_j})\mathbf{u_j} \\
                                   & = \sum_{k=1}^{l+m}(\mathbf{x}\cdot{} \mathbf{u_k})\mathbf{u_k}
                    \end{align*}

                    However, if the two subspaces aren't orthogonal complements, the union will not produce a orthonomal set.
                    This is because some of the unit vectors will be in the span of each other since the subspaces aren't orthogonal. \\

                    To fix this, all you have to do is apply the Gram–Schmidt algorthim on the union of the two old ONB to produce a new ONB for V-W.
          \end{enumerate}
    \item{} Intersection of two subspaces is also a subspace (nice little proof - you can do it later).
          \begin{enumerate}
              \item{} Argue that the intersection $A^{(1,2)} \cap A^{(5,6)}$ is either 0 ($zerosubspace$), 1 or 2 dimensional. Using the GS answer, argue it is $\{0\}$. \\

                    \textbf{SOLUTION:}

                    The subspace of the intersection has to be either 0, 1, or 2 dimensions because both subspaces are the span of two vectors:
                    \begin{align*}
                        A^{(1,2)} & = span(\mathbf{v_1}, \mathbf{v_2}) \\
                        A^{(5,6)} & = span(\mathbf{v_5}, \mathbf{v_6})
                    \end{align*}
                    There are only 3 possible outcomes:

                    1. Subspaces share no vectors, every vector is linearly independent, and there is no intersection, so the intersection is only the zero subspace because subspaces always pass through origin. \\
                    2. One vector lies in the span of the other subspace, therefore the dimension of the subspace of the intersection is 1, or a line of intersections. \\
                    3. Both pairs of vectors span the same subspace and the intersection has dimension two or a plane. \\

                    Using our vectors, during Gram–Schmidt, \(\mathbf{v_6}, \mathbf{v_5}\) both had orthogonal components to \(\mathbf{v_1}, \mathbf{v_2}\) because they did not produce the zero vector after subtracing off their projections to \(\mathbf{u_1},\mathbf{u_2}\) produced from \(\mathbf{v_1},\mathbf{v_2}\) respectively.
                    Therefore they are linearly independent and there is no intersection, so the subspace of the intersection is only the zero subspace.

              \item{} Argue that the intersection $A^{(1,2)} \cap A^{(4,5)}$ is either 0, 1 or 2 dimensional. Using the GS answer, argue it is $\{\alpha \mathbf{v}_4 : \alpha \in \mathbb{R}\}$. \\

                    \textbf{SOLUTION:}
                    Similarly to part a, the subspaces can be rewritten as:
                    \begin{align*}
                        A^{(1,2)} & = span(\mathbf{v_1}, \mathbf{v_2}) \\
                        A^{(4,5)} & = span(\mathbf{v_4}, \mathbf{v_5})
                    \end{align*}
                    Except we know from question 1(b), that \(\mathbf{v_4}\) is just a linear combiantion of {\(\mathbf{v_1}, \mathbf{v_2}\)}.
                    Therefore, the entire line that is parallel to \(\mathbf{v_4}\) is a intersection, the line of intersection can also be stated as:
                    \[
                        \{\alpha \mathbf{v}_4 : \alpha \in \mathbb{R}\}
                    \]

          \end{enumerate}
    \item{} Because they are two dimensional affine subsets, we consider $\mathbf{v}_7 + A^{(1,2)}$, $\mathbf{v}_6 + A^{(4,5)}$, and $A^{(5,6)}$ planes (just like what we visualize in $\mathbb{R}^3$) living in $\mathbb{R}^5$. Show that the first does not intersect with the second, and the first does not intersect with the last.

          \textbf{SOLUTION:}
          Back in question 8, we learned that the distance between two affine subsets can be written in a certain way:
          \begin{align*}
              \| \mathbf{p}-\mathbf{q} + V - W \| & = \|(\mathbf{p}-\mathbf{q}) - (V + W)\|
          \end{align*}
          For the affine subsets to intersect the distance between them must be zero,
          \begin{align*}
              \|(\mathbf{p}-\mathbf{q}) - (V + W)\| = 0 \\
              (\mathbf{p}-\mathbf{q}) - (V + W) = 0     \\
              (\mathbf{p}-\mathbf{q}) \in{} (V+W)
          \end{align*}
          Therefore, for the intersection to occur, the vector created from the subtracting one base point from another must lie inside the subspace formed by the combination of the two subspaces.

          From here rewrite the vectors as linear combinations of the ONB to see if they lie in the subspaces:
          \begin{align*}
              \mathbf{v}_7-\mathbf{v}_6 & \in{} A^{(4,5)} + A^{(1,2)}                                                                                                               \\
              \mathbf{v}_7-\mathbf{v}_6 & = \tfrac{1}{2}\mathbf{u_1} + \tfrac{\sqrt{3}}{3}\mathbf{u_2} - \tfrac{\sqrt{3}}{6}\mathbf{u_3} + \tfrac{\sqrt{3}}{3}\mathbf{u_5}-
              (\tfrac{1}{2}\mathbf{u}_1 + \tfrac{2\sqrt{3}}{3}\mathbf{u}_2 - \tfrac{\sqrt{3}}{6}\mathbf{u}_3 + \tfrac{\sqrt{3}}{3}\mathbf{u}_4)                                     \\
              A^{(4,5)} + A^{(1,2)}     & = s(2\mathbf{u_1}) + t(\sqrt{3}\mathbf{u_2})+ p(2\mathbf{u_1}+\sqrt{3}\mathbf{u_2}) + q(\tfrac{1}{2}\mathbf{u_1} + 2\sqrt{3}\mathbf{u_3})
              \quad{} | \quad{} (s,t,p,q \in{} \mathbb{R})
          \end{align*}
          It can be easily seen that $\mathbf{v}_7-\mathbf{v}_6$ contain orthogonal componets to the subspace spanned by $ A^{(4,5)} + A^{(1,2)}$, since both $\mathbf{u}_4, \mathbf{u_5}$ are orthogonal to the entire subspace.
          There is no way to minimize distance to zero, therefore there is no intersection here. \\

          A subspace is just a affine subset with basepoint 0:
          \begin{align*}
              \mathbf{v}_7-\mathbf{0} = \mathbf{v}_7 & \in{} A^{(5,6)} + A^{(1,2)} = span(\mathbf{v_1},\mathbf{v_2}, \mathbf{v_5})
          \end{align*}
          There is no reason to go through the extra work of expanding.
          From Gram–Schmidt we found that $\mathbf{v_7}$ revealed a new orthonormal vector \(\mathbf{u_5}\), therefore it can't possibly be in the span of the earlier vectors.
          The distance cannot be minimizied to zero.
          Again, there is no intersection.

    \item{} (a) Find the distance between the following sets
          \begin{enumerate}
              \item[(i)] $\{\mathbf{v}_7\}$ and $A^{(1,5)}$

                    \textbf{SOLUTION:}
                    To minimize this distance, first perform orthogonal decomposition on the distance between the sets:
                    \begin{align*}
                        \| \mathbf{v_7} - A^{(1,5)} \| & = \| (\mathbf{v_7} - A^{(1,5)})_{A^{(1,5)}} \| + \| (\mathbf{v_7} - A^{(1,5)}) - (\mathbf{v_7} - A^{(1,5)})_{A^{(1,5)}} \| \\
                                                       & =  \| \mathbf{v_7}_{A^{(1,5)}} - A^{(1,5)} \| + \| (\mathbf{v_7} - \mathbf{v_7}_{A^{(1,5)}} )\|
                    \end{align*}
                    The orthogonal component gives minimum distance since it is constant and we already know that \(A^{(1,5)} = span(\mathbf{u_1, u_2, u_3})\) from earlier:
                    \begin{align*}
                        \| (\mathbf{v_7} - \mathbf{v_7}_{A^{(1,5)}} )\| = \| \tfrac{1}{\sqrt{3}}\mathbf{u_3} \| = \tfrac{1}{\sqrt{3}}
                    \end{align*}
                    This is the minimized distance. \\

                    \emph{(b) Obtain the closest point(s), in each case.}

                    We can now find the closest point by setting the projection component to zero:
                    \begin{align*}
                        \mathbf{v_7}_{A^{(1,5)}} - A^{(1,5)}                                                         & =  0                                                      \\
                        \tfrac{1}{2}\mathbf{u}_1 + \tfrac{\sqrt{3}}{3}\mathbf{u}_2 - \tfrac{\sqrt{3}}{6}\mathbf{u}_3 & \in{} A^{(1,5)}                                           \\
                        (\tfrac{1}{3}, 0, \tfrac{1}{3}, 0, \tfrac{2}{3})                                             & = \mathbf{a} \quad{} | \quad{} \mathbf{a} \in{} A^{(1,5)}
                    \end{align*}
                    So \(\mathbf{a}\) is the closest point from the $A^{(1,5)}$ to $\mathbf{v_7}$.

              \item[(ii)] $\mathbf{v}_7 + A^{(1,2)} = \{\mathbf{v}_7 + r\mathbf{v}_1 + s\mathbf{v}_2 : r, s \in \mathbb{R}\}$ and $A^{(5,6)} = \{t\mathbf{v}_5 + u\mathbf{v}_6 : t, u \in \mathbb{R}\}$

                    \textbf{SOLUTION:}
                    From question 8(a), we have already derived a expression to minimze distance between two affine subsets.
                    \begin{align*}
                        \| \mathbf{p}-\mathbf{q} + V - W \| & = \| (\mathbf{p}-\mathbf{q})_{V-W} + (V - W)\| + \|(\mathbf{p}-\mathbf{q})-(\mathbf{p}-\mathbf{q})_{V-W} \|
                    \end{align*}
                    Plugging in our affine subsets we get:
                    \begin{align*}
                        \| \mathbf{v_7}-\mathbf{0} + A^{(1,2)} - A^{(5,6)} \| & =
                        \| (\mathbf{\mathbf{v_7}})_{A^{(1,2)}-A^{(5,6)}}  + A^{(1,2)} - A^{(5,6)}\| + \|\mathbf{v_7}-\mathbf{v_7}_{A^{(1,2)}-A^{(5,6)}} \|
                    \end{align*}
                    We can find the span of $(A^{(1,2)}-A^{(5,6)})$ in terms of the ONB using the equations we got in 1(b):
                    \begin{align*}
                        (A^{(1,2)}-A^{(5,6)}) & = span(\mathbf{v_1,v_2,v_5,v_6}) \\
                                              & = span(\mathbf{u_1,u_2,u_3,u_4})
                    \end{align*}
                    That means the only orthognal component of $\mathbf{v}_7$ is its $\mathbf{u_5}$ component:
                    \begin{align*}
                        \|\mathbf{v_7}-\mathbf{v_7}_{A^{(1,2)}-A^{(5,6)}} \| & = \| \tfrac{1}{\sqrt{3}}\mathbf{u_5} \| \\
                                                                             & = \tfrac{1}{\sqrt{3}}
                    \end{align*}
                    This is the minimized distance. \\

                    \emph{(b) Obtain the closest point(s), in each case.}

                    Because there is a closest point for each of the affine subsets, we need to solve for the parameters: r,s,t,u. \\
                    This can be done by solving the projection component as zero, minimizing it.
                    To do this we must first expand in terms of the ONB:
                    \begin{align*}
                         & \mathbf{\mathbf{v_7}}_{A^{(1,2)}-A^{(5,6)}} + A^{(1,2)} - A^{(5,6)} = 0                                       \\
                         & = \mathbf{\mathbf{v_7}}_{A^{(1,2)}-A^{(5,6)}} + r\mathbf{v_1} + s\mathbf{v_2} - t\mathbf{v_5} - u\mathbf{v_6} \\
                         & = (\tfrac{1}{2} + 2r - \tfrac{1}{2}t - \tfrac{1}{2}u) \mathbf{u_1}
                        + (\tfrac{\sqrt{3}}{3} + \sqrt{3}s - 2\tfrac{\sqrt{3}}{3}u) \mathbf{u_2}
                        + (-\tfrac{\sqrt{3}}{6} - \tfrac{\sqrt{3}}{2}t + \tfrac{\sqrt{3}}{6}u) \mathbf{u_3}
                        + (\tfrac{\sqrt{3}}{3}u) \mathbf{u_4}
                    \end{align*}
                    Then set each coefficient to zero, solving for the parameters:
                    \begin{align*}
                        \mathbf{u_4}: & \tfrac{\sqrt{3}}{3}u = 0                                                 \\
                                      & u = 0                                                                    \\
                        \mathbf{u_3}: & -\tfrac{\sqrt{3}}{6} - \tfrac{\sqrt{3}}{2}t + \tfrac{\sqrt{3}}{6}(0) = 0 \\
                                      & -\tfrac{\sqrt{3}}{6} = \tfrac{\sqrt{3}}{2}t                              \\
                                      & t = -\tfrac{1}{3}                                                        \\
                        \mathbf{u_2}: & \tfrac{\sqrt{3}}{3} + \sqrt{3}s - 2\tfrac{\sqrt{3}}{3}(0)=0              \\
                                      & \tfrac{\sqrt{3}}{3} = -\sqrt{3}s                                         \\
                                      & s = -\tfrac{1}{3}                                                        \\
                        \mathbf{u_1}: & \tfrac{1}{2} + 2r - \tfrac{1}{2}(-\tfrac{1}{3}) - \tfrac{1}{2}(0)=0      \\
                                      & 2r + \tfrac{4}{6} = 0                                                    \\
                                      & r = -\tfrac{1}{3}
                    \end{align*}
                    Plugging these parameters we have found the closest points for each subset,
                    \begin{align*}
                        \mathbf{p_1} & = \mathbf{v_7} + r\mathbf{v_1} + s\mathbf{v_2} \quad{} | \quad{} \mathbf{p_1} \in{} \mathbf{v_7} + A^{(1,2)} \\
                                     & =\mathbf{v_7} - \tfrac{1}{3} \mathbf{v_1} - \tfrac{1}{3} \mathbf{v_2}                                        \\
                                     & = \tfrac{1}{3} (-1,0,-1,-1,1)                                                                                \\
                        \mathbf{p_2} & = \mathbf{0} + t\mathbf{v_5} + u\mathbf{v_6} \quad{} | \quad{} \mathbf{p_2} \in{} A^{(5,6)}                  \\
                                     & = - \tfrac{1}{3} \mathbf{v_5} + 0\mathbf{v_6}                                                                \\
                                     & = \tfrac{1}{3} (0,0,0,-1,0)
                    \end{align*}

              \item[(iii)] $\mathbf{v}_7 + A^{(1,2)} = \{\mathbf{v}_7 + r\mathbf{v}_1 + s\mathbf{v}_2 : r, t \in \mathbb{R}\}$ and $\mathbf{v}_6 + A^{(4,5)} = \{\mathbf{v}_6 + t\mathbf{v}_4 + u\mathbf{v}_5 : t, u \in \mathbb{R}\}$

                    \textbf{SOLUTION:}
                    Use same strategy as the one in 11(ii):
                    \begin{align*}
                         & \| \mathbf{v_7} + A^{(1,2)} - \mathbf{v_6} - A^{(4,5)} \|                                                                                                              \\
                         & = \| (\mathbf{v_7 - v_6})_{A^{(1,2)} - A^{(4,5)}} + A^{(1,2)} - A^{(4,5)}\| + \|(\mathbf{v_7} - \mathbf{v_6}) - (\mathbf{v_7} - \mathbf{v_6})_{A^{(1,2)}-A^{(4,5)}} \| \\
                    \end{align*}
                    The minimum distance is the constant orthognal component.
                    We find the vector difference by subtracting the vectors in terms of the ONB,
                    then summing the orthognal components in regard to subspace \(A^{(1,2)} - A^{(4,5)} = span(\mathbf{v_1, v_2, v_4, v_5}) = span(\mathbf{u_1,u_2,u_3})\):
                    \begin{align*}
                        \mathbf{v_7-v_6}                                         & =
                        (\tfrac{1}{2}-\tfrac{1}{2}) \mathbf{u_1}
                        + (\tfrac{\sqrt{3}}{3}-\tfrac{2\sqrt{3}}{2}) \mathbf{u_2}
                        + (\tfrac{\sqrt{3}}{6}-\tfrac{2\sqrt{3}}{6}) \mathbf{u_3}
                        + (-\tfrac{\sqrt{3}}{3}) \mathbf{u_4}
                        + (\tfrac{\sqrt{3}}{3}) \mathbf{u_5}                                                                                                                                  \\
                                                                                 & =  -\tfrac{\sqrt{3}}{2} \mathbf{u_2} - \tfrac{\sqrt{3}}{4} \mathbf{u_4} + \tfrac{\sqrt{3}}{3} \mathbf{u_5} \\
                        \| (\mathbf{v_7-v_6})_{\perp{} A^{(1,2)} - A^{(4,5)}} \| & = \| - \tfrac{\sqrt{3}}{3} \mathbf{u_4} + \tfrac{\sqrt{3}}{3} \mathbf{u_5}          \|                     \\
                                                                                 & = \| \tfrac{1}{3}(0,1,-2,0,1)\|                                                                            \\
                                                                                 & = \tfrac{\sqrt{6}}{3}
                    \end{align*}
                    This is the minium distance. \\

                    From the last step, we have also found \(\mathbf{v_7-v_6}_{A^{(1,2)} - A^{(4,5)}} = -\tfrac{\sqrt{3}}{2}\mathbf{u_2}\). \\

                    \emph{(b) Obtain the closest point(s), in each case.}

                    We again use the expansion into ONB to discover the parameter values that will set the projected component to zero:
                    \\ (note that since $\mathbf{v_4}$ is in the span of $A^{(1,2)} \rightarrow{} A^{(4,5)} = A^5$)
                    \begin{align*}
                         & (\mathbf{v_7 - v_6})_{A^{(1,2)} - A^{(4,5)}} + A^{(1,2)} - A^{(4,5)} = 0 \\
                         & = (2r - \tfrac{1}{2}u) \mathbf{u_1}
                        + (\sqrt{3}s - \tfrac{2\sqrt{3}}{3}u - \tfrac{\sqrt{3}}{3}) \mathbf{u_2}
                        = (-\tfrac{3}{6}u) \mathbf{u_3}
                    \end{align*}
                    Set coefficients to zero and solve for parameters:
                    \begin{align*}
                        \mathbf{u_3}: & -\tfrac{3}{6}u = 0                                            \\
                                      & u = 0                                                         \\
                        \mathbf{u_2}: & \sqrt{3}s - \tfrac{2\sqrt{3}}{3}(0) - \tfrac{\sqrt{3}}{3} = 0 \\
                                      & \sqrt{3}s = \tfrac{\sqrt{3}}{3}                               \\
                                      & s = \tfrac{1}{3}                                              \\
                        \mathbf{u_1}: & 2r - \tfrac{1}{2}(0)=0                                        \\
                                      & r= 0
                    \end{align*}
                    Plugging these parameters we have found the closest points for each subset,
                    \begin{align*}
                        \mathbf{p_1} & = \mathbf{v_7} + r\mathbf{v_1} + s\mathbf{v_2} \quad{} | \quad{} \mathbf{p_1} \in{} \mathbf{v_7} + A^{(1,2)} \\
                                     & =\mathbf{v_7} + 0\mathbf{v_1} + \tfrac{1}{3} \mathbf{v_2}                                                    \\
                                     & = \tfrac{1}{3} (0,1,1,0,4)                                                                                   \\
                        \mathbf{p_2} & = \mathbf{v_6} + u\mathbf{v_5} \quad{} | \quad{} \mathbf{p_2} \in{} \mathbf{v_6} + A^{(4,5)}                 \\
                                     & = \mathbf{v_6} + 0\mathbf{v_5}                                                                               \\
                                     & = \mathbf{v_6}
                    \end{align*}

          \end{enumerate}
\end{enumerate}
\end{document}